\chapter{Validación profesional cualitativa}
\label{anexo:validacion-profesional}

Este anexo amplía la evidencia sintetizada en la sección~\ref{sec:resultados-validacion-profesional-cualitativa} del Capítulo~\ref{chap:resultados-evaluacion}. El cuerpo del Capítulo 13 presenta la síntesis; este anexo conserva el detalle de cada sesión, la trazabilidad hacia decisiones de producto y las limitaciones que acotan su alcance.

\section{Propósito y alcance}
\label{anexo:h-proposito-alcance}

Este trabajo distingue tres instancias de validación que no deben mezclarse entre sí:

\begin{itemize}
    \item \textbf{User research inicial (EN-01 a EN-05).} Entrevistas exploratorias documentadas en el Anexo~\ref{anexo:entrevistas} y en el Capítulo~8. Su objetivo fue descubrir necesidades y ajustar el alcance del prototipo antes de su desarrollo.
    \item \textbf{Validación profesional cualitativa (VP-01 y VP-02).} Las dos sesiones documentadas en este anexo. Su objetivo fue revisar una versión funcional/intermedia del producto ya construido, observando utilidad, flujo, comprensión, riesgos, confianza y prioridades de mejora.
    \item \textbf{Validación cruzada profesional final.} Todavía pendiente. Su objetivo futuro es que profesionales y el sistema analicen de forma independiente los mismos estudios, para luego revisar y documentar las discrepancias entre ambas lecturas (sección~\ref{anexo:h-validacion-cruzada-pendiente}).
\end{itemize}

Se utiliza la identificación VP-01 y VP-02 para las sesiones de este anexo, reservando la numeración EN- para las entrevistas exploratorias iniciales y evitando así que ambas instancias se confundan en la trazabilidad del proyecto.

Ninguna de las tres instancias, ni individual ni en conjunto, constituye validación clínica, validación diagnóstica, ensayo clínico, comparación contra ground truth ni evaluación de eficacia clínica.

\section{Metodología de las sesiones}
\label{anexo:h-metodologia}

Las sesiones VP-01 y VP-02 se organizaron como observación y comentario guiado de una versión funcional/intermedia del producto, utilizando como marco de referencia los criterios ya empleados en el proyecto para este tipo de instancia:

\begin{itemize}
    \item plausibilidad y claridad de la asistencia visual (overlays y máscaras);
    \item utilidad percibida de las mediciones ofrecidas;
    \item comprensión del flujo de trabajo propuesto;
    \item confianza en los resultados y fuentes concretas de desconfianza;
    \item posibilidad de corregir manualmente cada resultado;
    \item riesgo de sobreconfianza en la asistencia automática;
    \item funcionalidades consideradas indispensables por la profesional;
    \item utilidad general del producto como herramienta asistiva, no diagnóstica.
\end{itemize}

El material del proyecto incluye además una plantilla de referencia, \texttt{PF033\_professional\_\allowbreak validation\_\allowbreak template.md}, organizada en torno a los ejes de comprensión, utilidad, confianza, riesgo, flujo, primera versión y valor general. Esos ejes se retoman aquí como marco metodológico y de observación para estructurar las fichas de las secciones~\ref{anexo:h-vp01-torres} y~\ref{anexo:h-vp02-cejas}. No se afirma que cada profesional haya respondido literalmente cada uno de esos puntos ni que haya asignado puntuaciones numéricas: no existe evidencia documental de una escala de puntuación aplicada en estas sesiones, por lo que las fichas se presentan en formato narrativo y cualitativo, sin inventar valores cuantitativos, fechas, duración, matrícula, especialidad, institución ni citas textuales que no consten en la evidencia disponible.

\section{VP-01 --- Sesión con la Dra. Susana Torres}
\label{anexo:h-vp01-torres}

\subsection*{Profesional}
Dra. Susana Torres.

\subsection*{Contexto}
Sesión posterior a las entrevistas exploratorias iniciales (EN-01 a EN-05), metodológicamente distinta de ellas. La profesional observó y comentó una versión funcional/intermedia del prototipo, con foco en utilidad, flujo de revisión, visualización, confianza y priorización de mejoras.

\subsection*{Funciones mostradas o discutidas}
Segmentación como apoyo visual, overlays o máscaras sobre la imagen, posibilidad de apagar máscaras, corrección manual de mediciones, comparación longitudinal, correlación sagital--axial, revisión de múltiples cortes y salida revisable por el profesional.

\subsection*{Hallazgos documentados}
\begin{itemize}
    \item Interés en la comparación con resonancias previas y valor del historial longitudinal para observar la evolución de cambios degenerativos.
    \item La segmentación fue valorada como apoyo visual revisable, con preferencia por revisar primero la imagen sin asistencia y poder activarla después, para evitar sesgo en la lectura inicial.
    \item Posibilidad de apagar máscaras, corrección manual de mediciones y colores consistentes entre estructuras.
    \item Necesidad de evitar superposiciones que generen confusión visual.
    \item Relevancia del diámetro anteroposterior y transverso del saco o canal, y diferenciación de los neuroforámenes como entidad propia.
    \item Grasa alrededor de raíces y utilidad del producto para valorar estenosis.
    \item Importancia de revisar múltiples cortes en lugar de limitarse al corte central.
    \item Cambios degenerativos, hipertrofia del ligamento amarillo mencionada como posible extensión futura, y cambios Modic.
    \item Pérdida de altura discal.
    \item Correlación sagital--axial y respeto de la inclinación del disco en dicha correlación.
    \item Revisión profesional obligatoria como condición para confiar en el sistema.
\end{itemize}

Las sugerencias vinculadas a capacidades no implementadas (por ejemplo, hipertrofia del ligamento amarillo) se consignan como futuras y no deben interpretarse como funciones ya disponibles en el producto.

\subsection*{Problemas detectados}
Riesgo de sesgo si la máscara aparece antes de la lectura profesional; riesgo de confusión visual ante superposición excesiva; necesidad de evitar que la asistencia se interprete como diagnóstico; necesidad de no reducir la revisión a un único corte central.

\subsection*{Valoración general}
Favorable hacia un sistema de asistencia visual y documental, siempre que el profesional conserve la lectura, la corrección y la decisión final.

\subsection*{Limitaciones de esta sesión}
Sesión cualitativa, con una sola profesional, sobre una versión intermedia del producto. No incluye comparación contra lectura independiente, ground truth, sensibilidad ni especificidad, y no constituye validación clínica.

\section{VP-02 --- Sesión con la Dra. Claudia Cejas}
\label{anexo:h-vp02-cejas}

\subsection*{Profesional}
Dra. Claudia Cejas.

\subsection*{Contexto}
Sesión posterior a las entrevistas exploratorias iniciales (EN-01 a EN-05), orientada a evaluar utilidad y priorización funcional de una versión funcional/intermedia del prototipo.

\subsection*{Funciones mostradas o discutidas}
Identificación de nivel o contador, segmentación visual, corrección manual de mediciones, medición del canal, revisión de múltiples cortes, detección asistida de protrusiones focales o laterales, distinción entre canal central, canal lateral y neuroforamen, comparación longitudinal e informe estructurado editable.

\subsection*{Hallazgos documentados}
\begin{itemize}
    \item Utilidad del contador o identificación de nivel como apoyo práctico, y de la segmentación visual en general.
    \item Corrección manual de mediciones y prioridad del diámetro del canal.
    \item Recomendación de medir perpendicular al eje del canal.
    \item Necesidad de verificar bibliográficamente cualquier número o rango mencionado verbalmente antes de convertirlo en umbral del sistema.
    \item Importancia del neuroforamen, la estenosis foraminal y el receso subarticular o canal lateral.
    \item Necesidad de distinguir explícitamente canal central, canal lateral y neuroforamen.
    \item Revisión de múltiples cortes y detección de protrusiones focales o laterales.
    \item Cambios Modic, degeneración discal y pérdida de altura.
    \item Asociación entre el uso de colores/asistencia visual y la reducción de omisiones en contextos de carga laboral y fatiga.
    \item Cierta sobredetección considerada tolerable si el profesional conserva la posibilidad de revisar, corregir o descartar.
    \item Preferencia por detección y asistencia antes que diagnóstico automático.
    \item Utilidad potencial de la comparación longitudinal y valor de un informe estructurado editable.
\end{itemize}

\subsection*{Problemas detectados}
Riesgo de omisión de patología relevante; necesidad de revisar más de un corte; riesgo de convertir rangos mencionados verbalmente en umbrales sin respaldo bibliográfico; necesidad de mantener al profesional dentro del circuito de decisión.

\subsection*{Valoración general}
Favorable hacia una herramienta asistiva que reduzca omisiones y ordene el reporte, especialmente en contextos de carga laboral, siempre que no reemplace al profesional.

\subsection*{Limitaciones de esta sesión}
Sesión cualitativa, con una sola profesional, sobre una versión intermedia del producto. No constituye prueba clínica, comparación contra ground truth ni evaluación de sensibilidad o especificidad.

\section{Síntesis transversal de hallazgos}
\label{anexo:h-sintesis-transversal}

\subsection*{Usabilidad}
Las dos sesiones respaldan overlays útiles, máscaras activables o desactivables, edición manual, precisión visual y colores consistentes. El prototipo debe evitar superposiciones confusas y permitir que el profesional controle la asistencia visual.

\subsection*{Valor clínico-operativo}
Ambas profesionales valoraron la detección asistida, la reducción de omisiones, el apoyo ante carga laboral, el seguimiento longitudinal y la correlación sagital--axial. Esa valoración se interpreta como utilidad potencial de producto, no como validación clínica.

\subsection*{Mediciones prioritarias}
Se priorizan canal central, neuroforamen, receso lateral o subarticular y altura discal. Otras mediciones deben incorporarse solo si aportan valor real y pueden verificarse técnica o bibliográficamente.

\subsection*{Hallazgos relevantes}
Aparecen protrusión o herniación, estenosis, pérdida de altura, cambios degenerativos o Modic, forámenes y raíces. Algunos se encuentran fuera del producto validado y quedan como futuras líneas de trabajo.

\subsection*{Confianza y revisión humana}
La confianza depende de la revisión profesional obligatoria, la posibilidad de corregir, la transparencia sobre el estado de la IA y la ausencia de diagnóstico autónomo. Mostrar confianza puede ser útil cuando corresponda, pero no debe sustituir la revisión.

\subsection*{Diferencia de énfasis entre sesiones}
La Dra. Torres profundizó en comparación longitudinal, control visual de máscaras, correlación sagital--axial e inclinación discal. La Dra. Cejas enfatizó la protocolización de mediciones, la distinción entre canal central, canal lateral y neuroforamen, la carga laboral y la aceptabilidad de cierta sobredetección bajo revisión profesional.

\section{Trazabilidad entre observaciones y decisiones de producto}
\label{anexo:h-trazabilidad}

La Tabla~\ref{tab:h-cambios-producto-validacion-profesional} retoma la trazabilidad ya presentada en la Tabla~\ref{tab:res-cambios-producto-validacion-profesional} del Capítulo~13, entre cada observación profesional y la decisión de producto que motivó, junto con su estado actual. Los estados posibles son \texttt{IMPLEMENTED}, \texttt{PARTIAL}, \texttt{PLANNED} y \texttt{OUT\_OF\_SCOPE}; ningún estado se promueve sin evidencia técnica que lo respalde.

{\scriptsize
\setlength{\LTleft}{0pt}
\setlength{\LTright}{0pt}
\setlength{\tabcolsep}{3pt}
\renewcommand{\arraystretch}{1.18}
\begin{longtable}{>{\raggedright\arraybackslash}p{0.24\textwidth} >{\raggedright\arraybackslash}p{0.14\textwidth} >{\raggedright\arraybackslash}p{0.28\textwidth} >{\raggedright\arraybackslash}p{0.24\textwidth}}
\caption{Trazabilidad entre observaciones profesionales y decisiones de producto.}\label{tab:h-cambios-producto-validacion-profesional}\\
\hline
\textbf{Observación profesional} & \textbf{Profesional} & \textbf{Decisión de producto} & \textbf{Estado actual} \\
\hline
\endfirsthead
\hline
\textbf{Observación profesional} & \textbf{Profesional} & \textbf{Decisión de producto} & \textbf{Estado actual} \\
\hline
\endhead
\hline
\endfoot
Activar/desactivar máscaras y evitar confusión visual & Torres & Mantener overlays revisables y controles de visibilidad por estructura. & PARTIAL: overlays implementados; controles finos por estructura quedan como mejora de interfaz. \\
Corrección manual de mediciones & Torres/Cejas & Permitir que el profesional corrija valores y conserve trazabilidad. & IMPLEMENTED: mediciones editables y revisión profesional implementadas. \\
Evidencia visual verificable & Torres/Cejas & Mostrar máscaras, assets y regiones de evidencia asociadas al resultado. & IMPLEMENTED: overlays y evidencia visual servidos por Backend; limitados por registro AI efímero. \\
Comparación con estudios previos & Torres/Cejas & Mantener historial y comparación longitudinal descriptiva. & PARTIAL: implementado en Demo Mode/\texttt{subjectRef}; no es evaluación clínica de progresión. \\
Correlación sagital--axial y múltiples cortes & Torres/Cejas & Priorizar visor volumétrico, revisión por corte y no solo corte central. & PARTIAL: visor y catálogo full-series implementados; pairing sagital--axial 67C sigue experimental. \\
Medición de neuroforamen & Torres/Cejas & Registrar como medición prioritaria futura si se valida técnica y bibliográficamente. & PLANNED: no debe presentarse como medición productiva validada. \\
Receso lateral/subarticular & Cejas & Mantener línea subarticular como capacidad asistiva bajo gates. & PARTIAL: clasificador subarticular existe, pero P10.7/product path sigue bloqueado. \\
Transparencia Demo/IA y revisión humana & Torres/Cejas & Declarar modo demo/no clínico y mantener revisión obligatoria. & PARTIAL: flags técnicos existen; indicador persistente de Demo Mode en UI queda pendiente. \\
Informe estructurado editable & Cejas & Generar salida estructurada revisable y exportable. & IMPLEMENTED: reporte \texttt{pfi.study-report.v1}, PDF/JSON y revisión profesional. \\
Mostrar confianza cuando corresponda & Cejas & Comunicar confianza/estado solo cuando el modelo lo respalde. & PARTIAL: existen estados y gates; visualización sistemática de confianza queda pendiente. \\
Cambios Modic & Torres/Cejas & Documentar como hallazgo relevante, sin exponerlo productivamente sin validación. & PLANNED: P10.7 contiene tarea relacionada pero no está habilitada productivamente. \\
Hipertrofia del ligamento amarillo & Torres & Registrar como posible extensión futura. & PLANNED: no implementado en producto. \\
Diagnóstico automático & Torres/Cejas & Mantener enfoque asistivo, no diagnóstico. & OUT\_OF\_SCOPE: el sistema no diagnostica ni recomienda tratamiento. \\
\hline
\end{longtable}
}

\section{Limitaciones de esta instancia}
\label{anexo:h-limitaciones}

\begin{itemize}
    \item La muestra es de solo dos profesionales.
    \item La selección de profesionales fue intencional, no probabilística.
    \item La evaluación es cualitativa y se informa narrativamente, no mediante puntuaciones ni porcentajes.
    \item Ambas sesiones observaron una versión funcional/intermedia del producto, no la versión final consolidada.
    \item No se midió sensibilidad ni especificidad diagnóstica.
    \item No se comparó contra ground truth ni contra lectura profesional independiente sobre los mismos estudios.
    \item No se evaluó eficacia clínica.
    \item Los hallazgos no permiten generalizar a la totalidad de los profesionales potencialmente usuarios del producto.
    \item Algunas funcionalidades discutidas (por ejemplo, neuroforamen productivo, hipertrofia del ligamento amarillo) son parciales o futuras y no deben presentarse como implementadas por el solo hecho de haber sido propuestas o discutidas.
\end{itemize}

\section{Validación cruzada profesional final pendiente}
\label{anexo:h-validacion-cruzada-pendiente}

La validación cruzada profesional final continúa pendiente. El protocolo previsto para esa instancia futura consiste en:

\begin{enumerate}
    \item Seleccionar un mismo conjunto de estudios.
    \item Entregar esos estudios a profesionales para una lectura independiente, sin asistencia del sistema.
    \item Ejecutar los mismos estudios en el sistema.
    \item Registrar hallazgos y mediciones de ambos lados de forma separada.
    \item Comparar los resultados y documentar las discrepancias.
    \item Realizar una revisión posterior conjunta de esas discrepancias.
    \item Clasificar cada discrepancia según corresponda: coincidencia, omisión, sobredetección o diferencia de medición.
\end{enumerate}

Esta instancia no ha sido ejecutada al cierre de la presente versión, por lo que no se presentan resultados ni métricas de esa validación cruzada.
